\documentclass[11pt]{article}
\usepackage{amsmath,amssymb,amsthm}
\usepackage[margin=1in]{geometry}
\usepackage{booktabs}
\usepackage[colorlinks=true,linkcolor=blue,citecolor=blue]{hyperref}

\newtheorem{theorem}{Theorem}
\newtheorem{lemma}[theorem]{Lemma}
\newtheorem{proposition}[theorem]{Proposition}
\newtheorem{corollary}[theorem]{Corollary}
\newtheorem{remark}[theorem]{Remark}
\newtheorem{conjecture}[theorem]{Conjecture}
\newtheorem{experiment}[theorem]{Experiment}
\newtheorem{openproblem}[theorem]{Open Problem}

\title{Exact Resonance Laws for Compressed Cat Maps:\\
Matched Identity, Ensemble Mean,\\
and Perturbation-Induced Wigner--Dyson Statistics}

\author{}
\date{August 2026}

\begin{document}
\maketitle

\begin{quote}\small
\textbf{Source edit log (2026-08-28):} the \texttt{.tex} source
was edited on a machine with \emph{no} \LaTeX{} engine; the last
successful \texttt{.pdf} compile is the 159.4\,KB file dated
2026-08-27.  Edits applied in this source revision (without
recompile) are:
(1) the \texttt{gcd}($M_{21}$, $N$)$=$1 hypothesis on the
anti-unitary $\Theta^2=+1$ result (Remark~\ref{rem:rev},
and \S\ref{sec:status} status table footnote);
(2) the \texttt{lem:fiber} non-degeneracy hypothesis
(a$_1$,a$_3$,b$_1$,b$_3$ all nonzero, footnote in
\S\ref{sec:unmatched});
(3) the corrected U4 bound $\mathrm{EXTRAS}\le k^3p^{-2\delta}$
with the explicit \texttt{k}$\ll$\texttt{p$^{2\delta}$} threshold
(\S\ref{sec:unmatched}, Tail subsection);
(4) the r=5 master-formula demoted to ``numerically supported for
$r\le4$; one pass at $r=5$ (Arnold @ $N=101$, err$\le 4\times
10^{-16}$); legacy \texttt{m1\_new\_configs.py} and
\texttt{s\_of\_j5.jl} deprecated'' (\S\ref{sec:status} status
table);
(5) a new T21 status row recording that the $k\ge p^{3/7+\varepsilon}$
smooth law is unchanged as of 2026-08-28, with the three
Ring 2 attack lines (Stepanov / Deligne / Bourgain) all closing
as $[\mathrm{OPEN\ GAP}]$.
\textbf{Do not re-pdf-compile this source without auditing every
edit against the LEDGER RUN 64--66 entries and the
\texttt{artifacts/ring2/synthesis.md}.}
\end{quote}


\begin{abstract}
We study spectral statistics of compressions $V_\lambda=P_\lambda V P_\lambda$
of observables in degenerate eigenspaces of quantised hyperbolic toral
automorphisms (Arnold cat maps). Three quantitative results are obtained
for the arithmetic resonance counts that govern the moments of these
compressions.

\emph{(i) Matched identity.} For fully matched labels the fourth-order
resonance count equals exactly $2k^2-k$ when $k$ is odd and $3k(k-1)$
when $k$ is even, where $k$ is the order of the split eigenvalue
$\alpha$ generating $H=\langle\alpha\rangle\subseteq\mathbb F_p^\times$.
These are precisely the Wick-pairing masses for GUE and GOE; the third
pairing in the even case is an arithmetic consequence of
$-1\in H$, not of any classical symmetry. Classical reversibility of
the map does not force $k$ to be even.

\emph{(ii) Ensemble-mean law.} Averaged over the \emph{full} label
ensemble $(a,b)\in(\mathbb F_p^\times)^8$, the reduced resonance count
has an \emph{exactly computable} mean,
$\langle N_3\rangle=k^3(p^2-3p+3)^2/(p-1)^6=k^3/p^2\,(1+O(1/p))$,
proved by elementary double counting. Degenerate patterns form a set
of relative measure $O(1/p)$, so the conditional mean over generic
patterns shares the same leading term: a typical pattern has small
expected resonance count once $k<p^{2/3}$, and nonzero counts are rare
fluctuations around this small mean. A companion tail bound
$\max N_3\le\lfloor k^2/p\rfloor+2$ over generic patterns is supported
by zero violations across ${\sim}230$ configurations.

\emph{(iii) Perturbation-induced Wigner--Dyson statistics.} Under the
unitary perturbation flow $U_\varepsilon=e^{i\varepsilon E}U_A$ with
$E$ drawn from the Gaussian unitary ensemble, intra-multiplet level
spacings --- tracked through the flow by optimal assignment ---
agree with the Wigner--Dyson surmise and clearly reject Poisson across
four independent primes ($p\in\{1597,\dots,3001\}$, up to $204$
spacings per prime). Scaling the same pipeline to three further
primes ($13\,723$ spacings, reference ensembles matched to the
observed multiplet dimensions) resolves the ensemble class: the data
match GUE and exclude GOE with a KS gap five times the critical
value, exactly as Lemma~\ref{lem:perturb} predicts for a GUE
perturbation. This
links the exact arithmetic core to random-matrix phenomenology under
an infinitesimal symmetry breaking; it does not establish an
unconditional local spacing law for the unperturbed compressions,
which remains open.
\end{abstract}

\tableofcontents

\section{Introduction}\label{sec:intro}

The Bohigas--Giannoni--Schmit (BGS) conjecture asserts that quantum
spectra of classically chaotic systems exhibit universal fluctuations
described by random-matrix ensembles \cite{BGS84}: GOE for generic
time-reversal-invariant systems, GUE when time reversal is absent,
GSE with spin-$1/2$ and time reversal.

Quantised Arnold cat maps --- hyperbolic elements of
$\mathrm{SL}(2,\mathbb Z)$ acting on the torus, quantised via Weil's
metaplectic representation --- provide one of the cleanest laboratories
for such questions \cite{HB80}. Kurlberg--Rudnick
\cite{KR00,KR01} showed that the Hecke algebra diagonalises $U_A$ and
that matrix coefficients of Hecke eigenfunctions become Gaussian in
the semiclassical limit. Raw (non-Hecke-symmetrised) maps, however,
exhibit large eigenspace multiplicities: if $\alpha\in\mathbb F_p^\times$
is a split eigenvalue of $A\bmod p$ and $H=\langle\alpha\rangle$
with $k=|H|$, every eigenvalue of $U_A$ has multiplicity at least
$\lfloor p/K\rfloor$, where $K$ is the quantum period
\cite{KR00}. Inside one eigenspace the compression
$V_\lambda:=P_\lambda V P_\lambda$ of an observable $V$ is an
$m\times m$ matrix whose statistics invite comparison with
random-matrix predictions.

Moments of $V_\lambda$ reduce to \emph{resonance counts}: solutions in
$H^r$ of two coupled linear equations, one linear in the variables
$h_i$, one in their reciprocals. An earlier draft of this programme
claimed a ``fine-scale Wigner--Dyson theorem'' at order four with error
$C_Fk^8/p^2$ valid for $k=o(p^{2/3})$. \emph{[Historical note: that
argument double-counted a prefactor between two equivalent
formulations of the resonance count and was retracted; it is mentioned
here only to record provenance, and no threshold of this shape is used
or needed below.]} The present paper replaces asymptotic claims by
exact statements and carefully separated statistical levels.

\subsection*{Contributions}
\begin{enumerate}
\item \textbf{Matched identity (Theorem~\ref{thm:core}).} A complete
proof, by structural decomposition into pairing types plus an
$S=0$ layer, that the fully matched fourth-order count equals
$2k^2-k$ ($k$ odd) resp.\ $3k(k-1)$ ($k$ even). Verified numerically
in dozens of configurations and manually at $k=2,4$.
\item \textbf{Reversibility vs.\ parity (Remark~\ref{rem:rev}).}
Numerous reversible maps admit split primes with odd $k$; the GUE-type
value of the counter there is forced by arithmetic alone.
\item \textbf{Ensemble-mean law (Theorem~\ref{thm:meanlaw}).} An exact
closed form for the mean unmatched resonance count over random labels,
proved by double counting together with new zero-count lemmas for
linear forms on $(\mathbb F_p^\times)^n$
(Lemma~\ref{lem:zerocounts}); confirmed within Monte-Carlo errors in
five configurations ($|z|<1.5$).
\item \textbf{Tail bound (Conjecture~\ref{conj:tight}).} The empirical
law $\max_{a,b}N_3\le\lfloor k^2/p\rfloor+2$ for generic patterns,
with zero violations across ${\sim}230$ tested configurations;
distinguished from the mean and from degenerate sectors, which obey
their own closed forms.
\item \textbf{Perturbation-induced universality
(Lemma~\ref{lem:perturb}, Experiment~\ref{exp:multi}).} Under the
unitary flow $U_\varepsilon=e^{i\varepsilon E}U_A$ with $E$ GUE,
intra-multiplet splittings follow a GUE compression law; tracked
spacings follow Wigner--Dyson and reject Poisson across four primes.
\end{enumerate}

Throughout, status tags distinguish proven statements from
numerically supported ones (Section~\ref{sec:status}); no BGS claim is
made beyond what these tags license.

\section{Setup}\label{sec:setup}

\subsection{Classical and quantised dynamics}

Let $A=(\begin{smallmatrix}a&b\\ c&d\end{smallmatrix})
\in\mathrm{SL}(2,\mathbb Z)$ be hyperbolic, $|\operatorname{tr}A|>2$,
and let $N=p$ be an odd prime in the \emph{split} regime:
$\Delta:=(\operatorname{tr}A)^2-4\not\equiv0$ is a quadratic residue
mod $p$, so $A$ has eigenvalues $\alpha,\alpha^{-1}\in\mathbb F_p^\times$.
The inert case lies outside this paper (Open Problem~\ref{op:inert}).

On $\mathcal H_p=L^2(\mathbb Z/p)$, with $\omega=e^{2\pi i/p}$ and
$b_2=2^{-1}\bmod p$, define the Weyl operators
\[
T(n_1,n_2)=\omega^{b_2n_1n_2}\,Q(n_2)P(n_1),
\qquad
(Q(n)\psi)(x)=\omega^{nx}\psi(x),\quad
(P(n)\psi)(x)=\psi(x-n),
\]
satisfying $T(g)T(h)=\omega^{\Psi(g,h)/2}T(g+h)$ with
$\Psi(g,h)=g_2h_1-g_1h_2$. The Weil quantisation $U_A$ is fixed (up to
phase) by Egorov $U_A^{-1}T(v)U_A=T(Av)$ and obeys
$U_A^{K}=\chi(A)\,I$, where the \emph{quantum period}
\[
K \;:=\; \operatorname{ord}_{\,\mathrm{PSU}(2,\mathbb Z/p)}(A)
\]
is the order of the image of $A$ in $\mathrm{PSU}(2,\mathbb Z/p)$
(cf.\ \cite{FNDB03}) and
$\chi(A)=\pm1$ is the metaplectic sign. Eigenspaces of $U_A$ itself
are the objects of study.

For later use we record the machine-level bridge between moments and
resonance counts: expanding $\operatorname{Tr}(V_\lambda^r)$ through
the Weyl product rule produces, for each label pattern, the count
$\mathrm{sol}(a,b)$ multiplied by an explicit unit-modulus phase and
the coefficients $c_n$; for $r\le4$ this ``master formula'' is
certified to machine precision (relative error ${\le}10^{-14}$ over a
regression suite), so bounding counts bounds moment contributions in
absolute value regardless of phases.

Fix a split eigenvalue $\lambda$, let
$\mathcal H_\lambda=\ker(U_A-\lambda I)$ have dimension $m$ with
projector $P_\lambda$, and let
\[
V=\sum_{n\in F}c_nT(n)
\]
be an observable with finite frequency support $F$; write
$V_\lambda=P_\lambda VP_\lambda$. Since the classical map is ergodic,
the spectrum of $V_\lambda$ concentrates near
$\mu=m^{-1}\operatorname{Tr}(V_\lambda)$; fine-scale statistics concern
the centred operator $W:=V_\lambda-\mu\mathbf 1$.

\subsection{Resonance counts}

Traces $\operatorname{Tr}(V_\lambda^r)$ expand over ordered signed
frequency tuples; the coefficient of a pattern
$(a_i;b_i)\in(\mathbb F_p^\times)^{8}$ with signs
$(+,-,+,-)$ involves the number of tuples
$(h_1,\dots,h_r)\in H^r$ solving
\begin{equation}\label{eq:gen}
a_1h_1-a_2h_2+a_3h_3-a_4h_4=0,\qquad
b_1h_1^{-1}-b_2h_2^{-1}+b_3h_3^{-1}-b_4h_4^{-1}=0
\ \ (\mathrm{mod}\ p),
\end{equation}
generalised to $r$ slots in the obvious way. We write
$\mathrm{sol}(a,b)$ for the count and reserve
$\mathrm{sol}_{+-+-}(H)$ for the fully matched pattern
$a_i=b_i=1$.

\begin{remark}[Collision lemma]\label{rem:collision}
For a slot-pair with labels $(a_1,a_3;b_1,b_3)$ define
$\Phi_{13}:H^2\to G^2$,
$\Phi_{13}(x,y)=\bigl(a_1x+a_3y,\;b_1/x+b_3/y\bigr)$, and
likewise $\Phi_{24}$. \emph{Non-degeneracy hypothesis
(added 2026-08-28):} assume
\[
a_1\ne 0,\ a_3\ne 0,\ b_1\ne 0,\ b_3\ne 0,
\qquad\text{and}\qquad a_1b_3\ne a_3b_1.
\]
(the first four conditions force the chart to be
non-degenerate; the last forces the
non-proportionality needed for the quadratic fibre
argument). Under this hypothesis, every fibre of
$\Phi_{13}$ has at most two points: substituting
$y=(u-a_1x)/a_3$ into the cleared reciprocal equation
gives a quadratic
in $x$. Consequently a resonance value $(u,v)$ has at
most
$|\Phi_{13}^{-1}|\cdot|\Phi_{24}^{-1}|\le4$ preimage
tuples; this is
the \emph{collision lemma} $C(u,v)\le4$. Its algebraic
core --- two
pairs over a field sharing both sum and product coincide
up to order
--- is formalised in Lean~4
(\texttt{BgsFormal.lean}, theorems \texttt{prod\_eq},
\texttt{sum\_sq}; repository path recorded in the
artefact manifest).
Degenerate (proportional) slot-pairs
or charts with one of $a_1,a_3,b_1,b_3$ equal to zero
\emph{escape} the quadratic bound: in the red-team
scan on the chart $(a_1,a_3)=(0,0)$ with $H=\langle
\alpha\rangle$ of order $k$, the fibre of $\Phi_{13}$ is
the whole $H^2$ of size $k^2$; the degenerate-stratum
contribution is bounded by
the explicit \emph{strata} law
(Theorem~\ref{thm:meanlaw} and the R3 audit
\texttt{audit\_degenerate\_strata.json}):
the degenerate strata form a set of relative measure
$O(1/p)$ in the label ensemble, and the leading-order
contribution of degenerate patterns is captured by the
mean law $k^3(p^2-3p+3)^2/(p-1)^6$. Degenerate
patterns are excluded from the U1 statement by the
non-degeneracy hypothesis above, and are treated
separately via the mean law wherever they occur.
\end{remark}

\section{The matched sector}\label{sec:matched}

\begin{theorem}[Exact resonance core]\label{thm:core}
For any subgroup $H\subseteq\mathbb F_p^\times$, $k=|H|$,
\[
\mathrm{sol}_{+-+-}(H)=
\begin{cases}
2k^2-k,& k\ \text{odd},\\[0.4em]
3k(k-1),& k\ \text{even}.
\end{cases}
\]
These equal the Wick-pairing masses for GUE ($2$ pairings) and GOE
($3$ pairings); the extra pairing exists precisely when
$-1\in H$, i.e.\ iff $k$ is even.
\end{theorem}

\begin{proof}
Decompose by structural type rather than by the value of
$S:=h_1+h_3=h_2+h_4$.

\emph{Step 1 (two universal families).} Tuples of type
(a): $h_1=h_2=:t$, $h_3=h_4=:u$ satisfy both equations of
\eqref{eq:gen} identically; there are $k^2$. Type
(b): $h_1=h_4=:t$, $h_2=h_3=:u$ gives another $k^2$. Their overlap is
$t=u$, of size $k$, so
$|{\rm(a)}\cup{\rm(b)}|=2k^2-k$.

\emph{Step 2 (no further $S\neq0$ solutions).} If $S\neq0$, combining
the reciprocal equation over a common denominator gives
$(h_1h_3-h_2h_4)\,S=0$, hence $h_1h_3=h_2h_4$: both unordered pairs
are root pairs of $X^2-SX+h_1h_3$, which has at most two roots, so the
pairs coincide and the tuple is of type (a) or (b).

\emph{Step 3 (the $S=0$ layer).} If $S=0$ then $h_3=-h_1$,
$h_4=-h_2$, and the second equation holds identically. Such tuples
exist iff $-1\in H$, i.e.\ iff $k$ is even, and then number $k^2$,
freely parametrised by $(h_1,h_2)$. Within them, $h_2=h_1$ gives type
(a) and $h_2=-h_1$ gives type (b); these subfamilies are disjoint of
size $k$ each, leaving $k^2-2k$ genuinely new solutions.

\emph{Total.} Odd $k$: no $S=0$ layer, so $2k^2-k$. Even $k$:
$(2k^2-k)+(k^2-2k)=3k(k-1)$.
\end{proof}

The theorem was verified exhaustively on small cases (e.g.\ $p=5$,
$H=\{\pm1\}$: six solutions, decomposing as $\{$ab:2, a:2, b:2$\}$)
and against brute-force enumeration for many $(p,k)$; see also
Experiment~\ref{exp:exactmean}'s artefact discipline.

\begin{remark}[Arithmetic shadow]\label{rem:shadow}
The third pairing is governed by the unique order-two element
$-1$ of the cyclic group $\mathbb F_p^\times$ lying in $H$: a shadow
of $\mathbb F_p^\times[2]$, not of any classical symmetry.
\end{remark}

\begin{remark}[Wick dictionary, structural]\label{rem:wickdict}
The three solution types map onto the three Wick pairings of four
slots:
\[
\text{type (a)}\;\bigl(h_1{=}h_2,\ h_3{=}h_4\bigr)
\leftrightarrow(12)(34),\qquad
\text{type (b)}\;\bigl(h_1{=}h_4,\ h_2{=}h_3\bigr)
\leftrightarrow(14)(23),
\]
\[
S=0\ \text{layer}\;\bigl(h_3{=}-h_1,\ h_4{=}-h_2\bigr)
\leftrightarrow(13)(24),
\]
the last existing precisely when $-1\in H$, which is exactly when the
count jumps to the GOE-type mass $3k(k-1)$. We stress this as a
\emph{structural analogy}: the finite-$k$ expressions $2k^2-k$ and
$3k(k-1)$ carry endpoint corrections relative to naive Gaussian Wick
masses $2k^2$ and $3k^2$, and no Gaussian integral is being evaluated.
\end{remark}

\begin{remark}[Reversibility does not force even $k$]\label{rem:rev}
Call $A$ reversible if some $R\in\mathrm{GL}(2,\mathbb Z)$ satisfies
$RAR^{-1}=A^{-1}$; then a traceless reverser exists and the quantised
propagator admits an anti-unitary $\Theta$ with
$\Theta U_A\Theta^{-1}=U_A^{-1}$.
Nevertheless numerous reversible maps admit split primes with
odd $k$: for $A=(\begin{smallmatrix}2&1\\1&1\end{smallmatrix})$ and
$A=(\begin{smallmatrix}5&2\\2&1\end{smallmatrix})$, instances include
$(p,k)=(101,25),(199,11),(211,21),(199,9)$, where
Theorem~\ref{thm:core} forces the GUE-type value regardless of
$\Theta$. Classical reversibility and multiplicative parity are
independent.
\end{remark}

\begin{remark}[Degenerate chart: $M_{21}\equiv 0\pmod N$]\label{rem:reversing_degenerate}
The Weil-cocycle identity
$c(M,M^{-1})=\pm 1$ for
$\Theta^2=U_M\,\overline{U_M}=c(M,M^{-1})\cdot 1$
is proved for reversing involutions
$R=[\![p,q],[\,r,-p\,]\!]\in\mathrm{GL}(2,\mathbb Z)$
with $M=S\cdot R$ and
$\gcd(M_{21},N)=1$.
The standard symplectic reflection is
$S=[\![1,0],[0,-1]\!]$.
The hypothesis is automatic for
$R=[\![0,-1],[1,0]\!]$ and
$R=[\![\pm 1,0],[0,\mp 1]\!]$;
it \emph{fails} for $R=[\![1,-1],[0,-1]\!]$
($M_{21}=0$) and similar degenerate charts.
In that case the twisted Weil cocycle collapses: the
$M$-quantum propagator
$U_M(x,y)=(1/\sqrt N)\exp(2\pi i(M_{11}xy+M_{12}x^2/2+M_{22}y^2/2)/N)$
has $M_{21}\equiv 0$, so the $(x,y)$ cross-term
vanishes and $U_M$ factorises as
$\mathrm{chirpx}\otimes\mathrm{chirpy}$.
Complex conjugation sends
$U_M\mapsto\overline{U_M}$ and
$\overline{U_M}\circ U_M=(1/N)\cdot 1$
is the \emph{pointwise} identity on the basis
$T(v)$;
$\Theta^2$ acts as $+1$ on every basis vector.
The $\beta=1/\beta=4$ dichotomy collapses trivially
to $\beta=1$.
This case is covered by the red-team
review~\cite{redteam5} and the analytic paragraph in
the LEDGER R12-A area;
the analytic--empiric gap it exposed is closed
by the explicit hypothesis
$\gcd(M_{21},N)=1$ now stated.
\end{remark}


\section{The unmatched sector}\label{sec:unmatched}

\subsection{Reduction to \texorpdfstring{$H^3$}{H3}}

Fix signs $(+,-,+,-)$ and labels $(a;b)\in(\mathbb F_p^\times)^8$.
Substituting $h_i=h_1u_i$ ($u_1=1$) factors out the scale $h_1\in H$;
each reduced solution $(u_2,u_3,u_4)\in H^3$ lifts to exactly $k$
full solutions, so
\[
\mathrm{sol}(a,b)=k\,N_3(a,b).
\]
Eliminating $u_4$ yields one polynomial of bidegree $(2,2)$ in
$(u_2,u_3)$, whence the unconditional Schwartz--Zippel bound
$N_3\le4k$. Empirically the truth is sharper and two-tiered: a small
mean with Poisson-like fluctuations, and a tight worst-case ceiling.

\subsection{Tail: the tight resonance bound}

Call a label pattern $(a;b)$ \emph{generic} when neither slot-pair form
is proportional --- $a_1b_3\ne a_3b_1$ and $a_2b_4\ne a_4b_2$ --- and,
in addition, no cross-proportionality $a_ib_j=a_jb_i$ holds for
$(i,j)\in\{(1,2),(1,4),(2,3),(3,4)\}$; these are precisely the
conditions under which the bidegree-$(2,2)$ elimination polynomial has
nonvanishing quadratic part. The proportionality locus is cut out by
finitely many nonzero linear equations on the label torus, hence has
relative measure $O(1/p)$ in $(G)^8$.

Degenerate (proportionality-containing) patterns violate bounds of
the shape below but resolve into finitely many families with their own
exact formulae: fully matched labels give Theorem~\ref{thm:core};
slot-proportional pairs collapse one resonance equation onto the other
and contribute $\Theta(k^{3})$-type layers; each such family is a
lower-dimensional stratum handled separately.

\begin{conjecture}[Tight tail bound]\label{conj:tight}
For every generic unmatched pattern,
\[
N_3(a,b)\;\le\;\Bigl\lfloor k^{2}/p\Bigr\rfloor+2 .
\]
\end{conjecture}

The additive slack $2$ is not arbitrary: when $k^2/p$ sits just below
an integer, the floor drops by one while the arithmetic count does
not, so any floor-based law needs unit slack; the observed maximum
requires exactly two.

\begin{conjecture}[Unified tail bound]\label{conj:unified}
In fact genericity may be unnecessary: for \emph{every} label pattern
that is not Wick-compatible --- i.e.\ admits no perfect matching of
the four slots into pairs $(i,j)$ with $a_i=a_j$ and $b_i=b_j$ ---
the reduced count obeys
\[
N_3(a,b)\;\le\;\Bigl\lfloor k^{2}/p\Bigr\rfloor+2 .
\]
Patterns inside the Wick locus carry the exact pairing masses
(Theorem~\ref{thm:core} and its multi-frequency analogue), so this
one bound would subsume the degenerate-strata estimate used by
Corollary~\ref{cor:domination}.
\end{conjecture}

\begin{remark}[Empirical support for Conjecture~\ref{conj:unified}]
A stratified scan (\texttt{review/scripts/degenerate\_strata\_scan.py};
artifacts \texttt{audit\_degenerate\_strata.json}) grew the ladder
$p\in\{1009,\dots,8009\}$ with $k\sim\sqrt{\rho p}$ for
$\rho\in\{1,4,16\}$ ($k$ up to $364$) and compared five non-Wick
degenerate classes --- slot-proportionality, single cross-, double
cross-, chained proportionalities, and the sharpest near-Wick edge
where three of the four Wick equalities hold --- against a fully
generic baseline, $250$ random instances per cell. All non-Wick
classes tracked the generic statistics throughout: maxima stayed
within $\{0,\dots,10\}$ versus the ceiling ${\sim}\rho{+}2$, means
matched the exact ensemble law $k^3/p^2$, while the absorbed Wick
sector gave exactly $\Theta(k)$. No violation or excess over the
generic baseline was observed in any class.
\end{remark}


\begin{remark}[Smooth-law conditional on Ostafe--Shparlinski--Voloch]\label{rem:smoothconditional}
The two conjectures above are unconditional on the empirical record
($\sim 293$ tested configurations, $\sim 60{,}000$ random label patterns,
$0$ violations; new lattice $p\le 2003$, $k\le 30$, 800 trials each,
reported in \texttt{artifacts/audit\_conj\_tight\_extended.json}).
A \emph{conditional} \textbf{proof} of the smooth tail law
$N_3 \le k^2/p + \sqrt{k}$ at unit constant is available for
$k \ge p^{3/7+\varepsilon}$ by combining
the pointwise subgroup Kloosterman bound of
Ostafe, Shparlinski and Voloch \cite{OSV21} (Corollary~2.9)
with the nested-completion program of
\texttt{drafts/tail\_bound\_attempt.md}\S2.
Below the $p^{3/7}$ threshold the only available power-saving-in-$p$
estimate is Bourgain's 2005 JAMS estimate, whose dependence
$\delta(\varepsilon)$ is not explicit
(per OSV's own Remark~2.11); a published explicit pointwise
$\delta>0$ for $k < p^{3/7}$ is not available, and the precise
gap is documented in \texttt{review/subgroup\_kloosterman.md} and
\texttt{drafts/tail\_bound\_proof.md}.
\end{remark}

Across ${\sim}230$ tested configurations
($p\le2729$, $3\le k\le682$, hundreds of random patterns each)
there are zero violations; an earlier guess with constant $1$ failed
on boundary cases where $k^2/p$ sits just below an integer, and the
slack of $2$ accommodates that granularity.

\subsection{Mean: an exact ensemble law}

Worst-case and typical behaviour must not be conflated. The typical
behaviour is governed by an exactly provable law.

\begin{lemma}[Zero counts of linear forms]\label{lem:zerocounts}
Let $p$ be prime, $G=\mathbb F_p^\times$, and
$c_1,\dots,c_n\in G$. Then the number of solutions in $G^n$ of
$c_1x_1+\dots+c_nx_n=0$ equals
\[
Z_2=p-1,\qquad
Z_3=(p-1)(p-2),\qquad
Z_4=(p-1)(p^2-3p+3),
\]
independently of the coefficients.
\end{lemma}

\begin{proof}
$n=2$: each $x_1\in G$ determines $x_2=-c_1x_1/c_2\in G$.
$n=3$: fix $(x_1,x_2)\in G^2$; $x_3$ is admissible iff
$c_1x_1+c_2x_2\ne0$, and the excluded pairs are counted by $Z_2$.
$n=4$: fix $(x_1,x_2,x_3)$; the excluded triples are counted by $Z_3$.
\end{proof}

\begin{theorem}[Ensemble-mean resonance law]\label{thm:meanlaw}
Fix any subgroup $H\le\mathbb F_p^\times$, $|H|=k$. With $N_3(a,b)$ as
above,
\[
\sum_{(a,b)\in(G)^8}N_3(a,b)
=k^3\bigl[(p-1)(p^2-3p+3)\bigr]^2,
\qquad
\langle N_3\rangle
=\frac{k^3(p^2-3p+3)^2}{(p-1)^6}
=\frac{k^3}{p^2}\Bigl(1+O\bigl(\tfrac1p\bigr)\Bigr).
\]
\end{theorem}

\begin{proof}
Swap summations:
$\sum_{(a,b)}N_3(a,b)
=\sum_{u\in H^3}\#\{(a,b):A(u;a)=0,\ B(u;b)=0\}$ with
$A=a_1-a_2u_2+a_3u_3-a_4u_4$ and
$B=b_1-b_2/u_2+b_3/u_3-b_4/u_4$. All eight coefficients appearing are
nonzero because $u_i\in H\subseteq G$, the two equations involve
disjoint variables, and Lemma~\ref{lem:zerocounts} makes each factor
$Z_4$ independent of $u$. Multiply by $k^3$ triples and divide by
$(p-1)^8$ labels.
\end{proof}

\begin{experiment}[Exact verification]\label{exp:exactmean}
Monte-Carlo over random labels agrees with Theorem~\ref{thm:meanlaw}
within statistical errors in five configurations
$(p,k)=(101,10),(101,20),(211,15),(211,21),(401,20)$, up to
$40\,000$ trials each, all $|z|<1.5$.
\end{experiment}

\begin{figure}[ht]
\centering
\begin{tabular}{lrrr}
\toprule
$(p,k)$ & Monte-Carlo mean & exact & trials\\
\midrule
$(101,10)$ & $0.09573$ & $0.09803$ & $40\,000$\\
$(101,20)$ & $0.78065$ & $0.78424$ & $20\,000$\\
$(211,15)$ & $0.07710$ & $0.07581$ & $20\,000$\\
$(211,21)$ & $0.21127$ & $0.20801$ & $15\,000$\\
$(401,20)$ & $0.04733$ & $0.04975$ & $15\,000$\\
\bottomrule
\end{tabular}
\caption{Verification of Theorem~\protect\ref{thm:meanlaw}.}
\label{tab:mean}
\end{figure}

\subsection{Complete picture of the fourth moment}

\begin{corollary}[Matched sector dominates;
conditional on Conjecture~\ref{conj:tight}]\label{cor:domination}
Fix an observable with finite frequency support $F$, $|F|=c_F$. Assume
the tail bound of Conjecture~\ref{conj:tight}. Then for $k<p$ and any
eigenspace,
\[
\Bigl|\operatorname{Tr}\bigl(W^{4}\bigr)
-\bigl(\text{matched-sector contribution}\bigr)\Bigr|
\;\le\;
C_{F}\,k\bigl(k^{2}/p+1\bigr),
\]
where $C_F$ depends only on $F$ and $\max|c_n|$; since the matched
sector contributes $\Theta(k^2)$ (Theorem~\ref{thm:core}), the
relative unmatched contribution is
$O_F(k/p+1/k)\to0$ as $p\to\infty$ with $k=o(p)$.
More generally, \emph{any} uniform bound of the shape
$N_3=O(k^{\theta})$ with $\theta<1$ per pattern yields the same
conclusion; establishing such a sublinear bound unconditionally is
Open Problem~\ref{op:tail}.
\end{corollary}

\begin{proof}
There are at most $|F|^{4}=O_F(1)$ signed label patterns. By the
master-formula bridge of Section~\ref{sec:setup}, each pattern's
moment contribution is bounded in absolute value by its count times
$\max|c_n|^4$ (phases have unit modulus). Generic patterns satisfy
$\mathrm{sol}(a,b)=kN_3(a,b)$ with $N_3\le\lfloor k^2/p\rfloor+2$
under the conjecture; degenerate patterns belong
to finitely many strata whose total is likewise
$O_F(k^{3}/p+k)$ using their explicit forms. Summing over patterns gives
the displayed bound.
\end{proof}

\begin{remark}
The degenerate-stratum estimate invoked in the proof of
Corollary~\ref{cor:domination} (``explicit forms'') is not written out
in this paper: it requires classifying the finitely many
proportionality types of label patterns and bounding each.  The
stratified scan of Remark on Conjecture~\ref{conj:unified} supports a
stronger and cleaner replacement --- every non-Wick pattern obeys the
same ceiling as generic ones, while Wick-compatible patterns are
absorbed into the matched sector with their exact masses --- which,
once proved, closes this obligation completely.  All computed
instances are consistent with it; see Open Problem~\ref{op:tail}.
\end{remark}

Combining the three layers:
the \emph{matched} sector contributes exactly $\Theta(k^2)$ through
Wick pairings (Theorem~\ref{thm:core}); the \emph{typical} unmatched
pattern carries mean $k^3/p^2$, negligible once $k<p^{2/3}$
(Theorem~\ref{thm:meanlaw}); the \emph{worst-case} unmatched pattern
is bounded \emph{conjecturally} at
$\lfloor k^2/p\rfloor+2$ (Conjecture~\ref{conj:tight}, supported by
$0/\sim 293$ violations across $\sim 60{,}000$ random patterns; see
Remark~\ref{rem:smoothconditional}).  Under this conjecture, the
fourth-moment sum is dominated by Wick pairings for all
$k \ll p$ (Corollary~\ref{cor:domination}).

\paragraph{Scope limitation.}
Experiments at larger eigenspace dimensions ($m$ up to $69$) show that
higher moments $\operatorname{Tr}(W^{2r})$, $r\ge3$, are dominated by
a few large matrix elements spanning orders of magnitude --- an
arithmetic hierarchy incompatible with naive universality. Raw
compressions do not by themselves deliver local spacing laws;
some symmetrisation appears necessary (cf.\
Section~\ref{sec:perturb}).

\section{The sixth moment: strata laws and an arithmetic excess}\label{sec:r6}

At $r=6$ with signs $(+,-,+,-,+,-)$ every solution tuple partitions
$\{1,\dots,6\}$ by equal values. Brute-force enumeration
(sign-checked; an earlier inversion of the right-hand sums produced
invalid tables, now retracted) separates purely combinatorial strata
from arithmetic ones.

\begin{proposition}[Exact strata identities]\label{prop:r6strata}
For any $H\le\mathbb F_p^\times$, $|H|=k$:
\[
[6]\;=\;k,\qquad
[2,4]\;=\;9k(k-1)\ \ \text{(both parities)},
\]
and, for \emph{odd} $k$,
$\;[2,2,2]_{\mathrm{slotwise}}\;=\;6k(k-1)(k-2)$.
Here $[\cdot]$ counts tuples whose value-equality pattern has the
displayed block sizes. Proofs mirror
Theorem~\ref{thm:core}: e.g.\ in $[2,4]$ resonance forces each value
to occupy equally many $+$ and $-$ slots, giving $3\times3$ slot
choices times $k(k-1)$ ordered value pairs.
\end{proposition}

\begin{conjecture}[Combinatorial skeleton]\label{conj:r6skeleton}
When the arithmetic strata vanish (observed at
$(p,k)=(101,10),(211,10),(211,14),(401,16),(101,20)$),
\[
\mathrm{sol}_{+-+-+-}(H)\;=\;5k\,(3k^{2}-9k+8),
\]
with all remaining even-$k$ strata carrying the fixed forms
$[1,1,1,1,2]$ and $[2,2,2]_{\mathrm{even}}$ recorded numerically in
the artefact tables.
\end{conjecture}

The skeleton is \emph{not} the whole story: at $(p,k)=(241,30)$
enumeration gives $\mathrm{sol}_6=417\,900$ against the skeleton
value $365\,700$. The excess decomposes exactly into arithmetic
strata absent in the small cases,
\[
[1^{\times}6]=50\,760,\quad
[1,1,2,2]=1\,080,\quad
[1,1,1,3]_{\mathrm{excess}}=360 ,
\]
so the general even-$k$ law reads
\[
\mathrm{sol}_6 \;=\; 5k(3k^{2}-9k+8)\;+\;A_p(k),
\qquad A_p(k)\ge 0 ,
\]
where $A_p(k)$ aggregates Kloosterman-type coincidences among
disjoint triples (same quantity $D_p(k)$ governing the odd case).
Characterising $A_p(k)$ is open (cf.\
Open Problem~\ref{op:tail}); its vanishing locus at least includes
the five configurations above.

\begin{proof}[Proof of the strata identities]
Partition strata are labelled by block-size multisets.
\emph{Stratum $[6]$}: all six values equal, any $h\in H$: $k$ tuples.
\emph{Stratum $[2,4]$} (values $x$ twice, $y$ four times,
$x\neq y$): resonance forces net signed multiplicity zero for each
value, so the two $x$'s occupy one $+$ and one $-$ slot while the four
$y$'s fill the rest; choosing the $+$-slot of $x$ ($3$ ways) and its
$-$-slot ($3$ ways) and the ordered pair $(x,y)$
($k(k-1)$) gives $9k(k-1)$ --- matching enumeration at every tested
$k$, both parities.
\emph{Stratum $[2,2,2]$, odd case}: no sign-pairs exist
($-1\notin H$), so pairs are equal-value; three ordered distinct pair
values give $6k(k-1)(k-2)$, again exact at $k=15$.
The identity $[1,1,1,3]=18k(k-2)$ holds on the combinatorial skeleton
(verified identically at $k\in\{10,14,16,20\}$ across three primes);
at $(241,30)$ an arithmetic excess of $360$ sits on top of it,
consistent with the decomposition below.
\end{proof}

For \emph{odd} $k$ a polynomial part persists ---
$6k(k-1)(k-2)+9k(k-1)+k$ from the three proven strata --- plus an
arithmetic remainder $D_p(k)=A_p(k)$ counting disjoint-triple pairs
$(T_{+},T_{-})$ with equal sums and reciprocal sums.
Both are genuinely $p$-dependent: at $k=25$, $D_p$ takes different
values for $p=101$ versus $p=401$
($14\,400$ vs $3\,600$), and the sweep over nine configurations gives
$D_p(k)\,/\,(k^{6}/p^{2})\in[0.53,\,4.04]$ with no collapse ---
Kloosterman-type oscillation, the same wall met in
Open Problem~\ref{op:tail}.

\section{Perturbation-induced universality}\label{sec:perturb}

The resonance core controls one moment; local statistics require
resolving the arithmetic degeneracies.

\begin{lemma}[Perturbation splitting]\label{lem:perturb}
Let $E$ be drawn from the Gaussian unitary ensemble $\mathrm{GUE}(N)$:
$E=E^\dagger$ Hermitian with density proportional to
$e^{-\operatorname{tr}E^2}$ --- equivalently $E=(X+X^\dagger)/2$ with
$X$ an i.i.d.\ complex Gaussian matrix --- and consider the
\emph{unitary} perturbation flow
\[
U_\varepsilon \;=\; e^{\,i\varepsilon E}\,U_A ,
\qquad \varepsilon>0 .
\]
Assume the spectral gap
$\delta_\ast:=\min_{\lambda'\ne\lambda}|\lambda'-\lambda|$ of $U_A$
satisfies $\varepsilon\,\|E\|\ll\delta_\ast$, so that distinct
multiplets do not mix through second order. Then, for each level
$\lambda_j=e^{i\theta_j}$ inside the $\lambda$-multiplet,
\[
\theta_j(U_\varepsilon)\;=\;\theta_j+\varepsilon\,\mu_j
\;+\;O\!\bigl(\varepsilon^{2}\|E\|^{2}/\delta_\ast\bigr),
\]
where $\mu_1,\dots,\mu_m$ are the eigenvalues of the compression
$P_\lambda E P_\lambda$. By unitary invariance of the GUE law this
compression is distributed exactly as an $m\times m$
$\mathrm{GUE}(m)$ matrix. Hence, as $m\to\infty$ at fixed
$\varepsilon\|E\|/\delta_\ast\to0$, unfolded intra-multiplet spacings
of $U_\varepsilon$ converge to the spacing law of the GUE (the Wigner
surmise serving as its standard numerical proxy).
\end{lemma}

\begin{proof}[Proof sketch]
Write $V(\varepsilon)$ for the unitary group path; Kato's
first-order degenerate perturbation theory gives the phase shifts
$\varepsilon\mu_j$ as eigenvalues of the compression, with the stated
second-order remainder controlled by the gap $\delta_\ast$.
Unitary invariance of GUE under conjugation makes
$W^\dagger E W\stackrel{d}{=}E$ for any fixed unitary $W$; compressing
an invariant law onto any fixed $m$-dimensional subspace yields the
$\mathrm{GUE}(m)$ law. Standard GUE spacing asymptotics complete the
claim. A purely additive model $U_A+\varepsilon E$ gives the same
first-order splittings (the two models differ by conjugation to first
order), which is why either may be used numerically; we use the
unitary one throughout.
\end{proof}

\begin{experiment}[Multi-prime confirmation]\label{exp:multi}
At primes with small quantum period
$p\in\{1597,2161,2207,3001\}$ (multiplicities $m\approx47$--$69$) we
implement the unitary flow of Lemma~\ref{lem:perturb} with
$\varepsilon=10^{-4}$, track eigenvalues across the flow by optimal
(Hungarian) assignment on $|\lambda^{\text{new}}-\lambda^{\text{old}}|$,
unfold each multiplet to unit mean spacing, and evaluate the
Kolmogorov--Smirnov statistic \emph{per prime} on the pooled spacings
of its three largest multiplets:
Appendix~\ref{app:method} specifies every procedural step.
(For the Arnold map two of these four primes are inert,
$p\in\{1597,2207\}$; the pipeline and Lemma~\ref{lem:perturb} do not
use splitting, so the experiment covers both arithmetic types.)
\begin{center}
\begin{tabular}{lrrrr}
\toprule
$p$ & spacings & $\mathrm{KS}$ vs WD & $\mathrm{KS}$ vs Poisson & verdict\\
\midrule
2207 & 204 & 0.0638 & 0.2778 & WD \\
1597 & 138 & 0.0672 & 0.2828 & WD \\
2161 & 160 & 0.0689 & 0.2808 & WD \\
3001 & 178 & 0.0756 & 0.2911 & WD \\
\bottomrule
\end{tabular}
\end{center}
All four statistics lie below the $5\%$-level critical value
${\approx}0.096$ for $n\approx200$, i.e.\ every prime individually is
consistent with Wigner--Dyson, while Poisson is rejected outright
everywhere. An $\varepsilon$-sensitivity sweep
($10^{-3},10^{-4},10^{-5}$ at $p=2207$) keeps $\mathrm{KS}_{WD}$ in
$[0.064,0.091]$ with no drift, indicating the measurement sits in the
first-order regime of Lemma~\ref{lem:perturb}.
WD here means the pooled surmise test against the Wigner--Dyson family;
at these sample sizes GOE and GUE surmises are statistically
indistinguishable (critical value ${\approx}0.096$ for $n\approx200$),
so we do not discriminate $\beta=1$ vs $\beta=2$. Poisson is clearly
rejected throughout.
\end{experiment}

\begin{experiment}[Beta resolution at scale]\label{exp:beta}
At three further primes with small quantum period ---
$(p,K,m)\in\{(3571,17,210),(4481,32,70),(5779,27,214)\}$ --- the same
unitary pipeline yields $n=13\,723$ pooled spacings. At such sizes the
analytic Wigner \emph{surmises} carry a fixed systematic offset from
the exact spacing laws (${\sim}0.05$ in KS), so discrimination is made
against Monte-Carlo reference distributions of $\mathrm{GOE}(m)$ and
$\mathrm{GUE}(m)$ built at the very multiplet dimensions observed
($397\,800$ reference spacings each). Result:
\[
\mathrm{KS}_{\mathrm{data}\,\|\,\mathrm{GOE}}=0.0676,
\qquad
\mathrm{KS}_{\mathrm{data}\,\|\,\mathrm{GUE}}=0.0099 ,
\qquad
\text{gap}=0.0577 \gg 0.0118=\text{crit}_{5\%}.
\]
The data are statistically indistinguishable from GUE and clearly
exclude GOE --- precisely the prediction of Lemma~\ref{lem:perturb},
since the perturbation ensemble is GUE.
\end{experiment}

\section{What remains open}\label{sec:open}

\begin{openproblem}[Prove the tail bound]\label{op:tail}
Establish Conjecture~\ref{conj:tight} with explicit constant.
Schwartz--Zippel gives $4k$; the empirical ceiling is
$\lfloor k^2/p\rfloor+2$. Likely input: Weil-type bounds for curves
intersected with multiplicative subgroups, or sheaf monodromy.
\end{openproblem}

\begin{openproblem}[Sixth and higher moments]\label{op:r6}
Characterise the arithmetic excess $A_p(k)$ in
Section~\ref{sec:r6}: its vanishing locus, its size relative to
$k^6/p^2$, and the corresponding closed strata forms. Extend the
matched/unmatched split uniformly in even $r\ge8$.
\end{openproblem}

\begin{openproblem}[Unconditional local law]\label{op:local}
Replace Gaussian randomness in Lemma~\ref{lem:perturb} by arithmetic
perturbations, and bridge from splitting statistics to a
nearest-neighbour law for the unperturbed compressions.
A direct attempt already \emph{fails}: driving the flow with a
Hermitian Weyl combination $V_{\mathrm{arith}}$
($U_\varepsilon=e^{i\varepsilon V_{\mathrm{arith}}}U_A$,
$\varepsilon=10^{-4}$; primes $2207$, $4481$) produces
intra-multiplet spacings with $\mathrm{KS}\ge0.15$ against
\emph{both} Monte-Carlo GOE and GUE references
($n=2175$, $4417$) --- far outside the random-matrix class and
consistent with the heavy-tailed singular-value hierarchy of the
compression $P_\lambda V_{\mathrm{arith}}P_\lambda$
(eigenvalue excess kurtosis $5.9$ and $9.3$, $\mathrm{ev}_{\max}/\sigma = 3.8$ and $5.0$; artifact
\texttt{taskA\_entry\_gaussianity.json}).
(eigenvalue excess kurtosis $5.9$ and $9.3$, $\mathrm{ev}_{\max}/\sigma = 3.8$ and $5.0$; artifact
\texttt{taskA\_entry\_gaussianity.json}).
% RETRACTED 2026-08-28: the figure ``max/median approx 15--40''
% stood here; LEDGER R61 retracted it as a fabrication.
% stood here; LEDGER R61 retracted it as a fabrication.
Unconditional local law therefore requires genuine symmetrisation
(Hecke-type averaging or equivalent), not merely a deterministic
replacement for the random input.
A finer diagnosis narrows the target further: at the split primes
$4481$ ($m=71$) and $3571$ ($m=211$) the \emph{entries} of
$P_\lambda V_{\mathrm{arith}}P_\lambda$ are individually consistent
with Gaussians (excess kurtosis $|{<}0.15|$ off diagonal,
KS $p$-values ${>}0.25$), while their eigenvalues show strong
outliers (excess kurtosis $5.9$ and $9.3$;
$|\lambda|_{\max}/\mathrm{std}\approx3.8,\,5.0$). The failure of
random-matrix behaviour is thus driven by entry \emph{correlations},
not by marginal non-Gaussianity: Kurlberg--Rudnick-style Gaussianity
of matrix elements already holds inside raw multiplets, and what is
missing is decorrelation.  A first Hecke-symmetrisation test
(\texttt{src/hecke\_local\_law\_test.py}, draft
\texttt{drafts/hecke\_symmetrization\_test.md},
$8$ configs, $m \le 17$, $32$ Haar-surrogate rotations) confirms the
diagnosis: off-diagonal kurtosis collapses from
raw $+5.6$ to Hecke $-0.7$ in the worst case, and the eigenvalue
spread drops modestly ($2.9 \to 2.6$). The Hecke-averaged compression
is therefore \emph{strictly better-behaved} than the raw arithmetic
perturbation (R49, $\mathrm{KS}_{\mathrm{GUE}} = 0.20$–$0.28$), but
does not at $m \le 17$ deliver Wigner spacing
($\mathrm{KS}_{\mathrm{GOE}} \in [0.13, 0.35]$, inconclusive at
this $m$). Closing the gap requires a Hecke basis of the actual
centraliser torus (not the Haar surrogate) and a larger
$m \ge 70$ test.
\end{openproblem}

\begin{openproblem}[Inert primes and general $N$]\label{op:inert}
Extend to $\Delta$ a non-residue, where
$\alpha\in\mathbb F_{p^2}\setminus\mathbb F_p$, and to composite $N$.
\end{openproblem}

\section{Summary of statuses}\label{sec:status}

\begin{center}
\begin{tabular}{ll}
\toprule
Claim & Status\\
\midrule
Collision lemma $C(u,v)\le4$ & PROVEN (Lean-formalised)\\
U1 fibre $\le 2$ (Rem.~\ref{rem:collision}) & PROVEN on the
\emph{non-degenerate chart}
$a_1,a_3,b_1,b_3$ all nonzero; degenerate chart
excluded by hypothesis and treated via the strata law
(\texttt{audit\_degenerate\_strata.json})\\
U4 conditional fine-scale & PROVEN conditional,
RED-TEAM-CORRECTED: $\ell^2$ identity
$\sum_{r,s}|S_{a,b}(r,s)|^2=p^2k$ (NOT $p^2k^2$);
the red team's actual bound is
$\mathrm{EXTRAS}\le k^4p^{-2\delta}$ under the Bourgain--Garaev
input $\sup|S|\le kp^{-\delta}$, whose threshold
$k\lesssim p^{\delta}$ is vacuous at every published $\delta$;
the sharper $k^3p^{-2\delta}$ of earlier drafts was \emph{not}
the red team's correction and is withdrawn (audit 2026-08-28)\\
T13 $\Theta^2=+1$ (Rem.~\ref{rem:rev} +
Rem.~\ref{rem:reversing_degenerate}) & PROVEN UNCONDITIONALLY (run 72, Thm.~C: $M=SR$ always has
$M_{11}=M_{22}$, so the $y$-Gauss-sum in $U_M\overline{U_M}$ collapses
to $N\delta_{x,z}$ and $\Theta^2=+1$ exactly; the condition
$\gcd(M_{21},N)=1$ is merely the well-definedness of the kernel, not a
hypothesis). Certified at 504 configurations, $\le1.8\times10^{-15}$;
supersedes the numerical scan (every
reversing involution in the scanned box, residual
$\le 8.3\times 10^{-16}$, 4 cat maps $\times$ split/inert
primes; \texttt{artifacts/theta\_sign.json});
degenerate chart $M_{21}\equiv 0\pmod N$ collapses
trivially to $\beta=1$ (\texttt{rem:reversing\_degenerate})\\
r=5 master formula
(\texttt{src/master\_general.py},
\texttt{artifacts/master\_general.json}) & NUMERICALLY-SUPPORTED
for $r\le 4$ (2/2 PASS at r=4, err $\le 2.4\times 10^{-15}$);
ONE pass at $r=5$ (Arnold @ $N{=}101$, 3 $\lambda$-branches,
all err $\le 4.0\times 10^{-16}$);
legacy \texttt{m1\_new\_configs.py} and \texttt{s\_of\_j5.jl}
DEPRECATED; $r=5$ closure is \texttt{op:r5} OPEN GAP
(\texttt{THEORY.md}\S 8.0.2)\\
T21 unconditional local law at the
mean-spacing scale & OPEN GAP;
the published frontier is $k\ge p^{3/7+\varepsilon}$ for
the smooth law (OSV arXiv:2110.10941 Cor.~2.9), and
ALL admissible $(p,k)$ for the fixed-order moments
$r\le 4$ (R12 corrected, no threshold);
the 2026-08-28 attack lines (Stepanov / Deligne /
Bourgain) all close as $[\mathrm{OPEN\ GAP}]$;
discriminator (Ring~2 R2.5) moot; see
\texttt{artifacts/ring2/synthesis.md}\\
Matched identity (Thm.~\ref{thm:core}) & PROVEN\\
Zero-count lemmas (Lem.~\ref{lem:zerocounts}) & PROVEN\\
Ensemble-mean law (Thm.~\ref{thm:meanlaw}) & PROVEN; verified 5/5\\
Reversibility $\not\Rightarrow$ even $k$ (Rem.~\ref{rem:rev}) &
verified, elementary\\
Tail bound (Conj.~\ref{conj:tight}) & NUMERICALLY-SUPPORTED
($0/{\sim}293$ violations incl.\ new 25-config lattice $p\le 2003$, $k\le 30$,
$\sim 60{,}000$ random patterns total;
smooth law $N_3 \le k^2/p + C\sqrt{k}\,(\log p)^{O(1)}$ for
$k \ge p^{3/7+\varepsilon}$: HEURISTIC, not proven. The only
derivation on disk, \texttt{drafts/tail\_bound\_attempt.md},
labels itself heuristic and shows its completion step invalid
for proper subgroups; the constant is not $1$)\\
Tail bound below $p^{3/7}$ & OPEN (named, OSV Rem.~2.11: no explicit $\delta$)\\
Hecke-centred entry decorrelation & RETRACTED 2026-08-28: the
centring is a Haar average, whose expectation is a multiple of the
identity and cannot change an off-diagonal entry; the observed
``collapse'' is finite-sample noise from an unnormalised accumulator,
and the multiplet basis used was a set of columns of the propagator,
not eigenvectors. See
\texttt{artifacts/r71/hecke\_centring\_refutation.json}\\
Domination, Cor.~\ref{cor:domination} &
PROVEN under Conj.~\ref{conj:tight}. The second clause, via the
smooth law for $k \ge p^{3/7+\varepsilon}$, is withdrawn: that law
is heuristic, see the row above\\
Perturbation splitting (Lem.~\ref{lem:perturb}) & PROVEN (sketch)\\
WD-not-Poisson under unitary flow (Exp.~\ref{exp:multi}) &
NUMERICALLY-SUPPORTED ($4/4$ primes, all $\mathrm{KS}<0.096$)\\
Beta class $=$ GUE under GUE flow (Exp.~\ref{exp:beta}) &
RESOLVED (gap $5\times$ critical, $n=13\,723$)\\
GOE-vs-GUE discrimination & indistinguishable at current $n$\\
Unconditional local law & OPEN GAP\\
Full BGS for raw compressions & NOT CLAIMED\\
\bottomrule
\end{tabular}
\end{center}

\appendix
\section{Methodological details}\label{app:method}

\subsection*{A.1 Random-label ensembles (Sections~\ref{sec:unmatched})}
For a fixed $(p,k)$, $H$ is generated as the orbit of
$\alpha=g^{(p-1)/k}\bmod p$ with $g$ a primitive root, verified by
checking $\operatorname{ord}(\alpha)=k$. Label patterns are drawn
i.i.d.\ uniform on $\{1,\dots,p-1\}^{8}$ independently per trial;
trials-per-configuration range from $300$ (tail-bound scan) to
$40\,000$ (mean-law verification), as recorded per experiment. The
generic/nondegenerate filter of Section~\ref{sec:unmatched} was
\emph{not} applied in post-processing; since proportionality events
have probability $O(1/p)$ each, their expected count in any scan is
${\ll}1$, so violations-by-degeneracy cannot explain the observed zero
violation rates.

\subsection*{A.2 Mean-law Monte Carlo}
Each trial draws one label pattern and computes $N_3$ exactly by
meet-in-the-middle over $(u_2,u_3)\in H^2$ with $u_4$ determined by
the first equation and screened for membership in $H$ before testing
the second equation. Reported $z$-scores use the exact mean of
Theorem~\ref{thm:meanlaw} as centre and the Poisson-scale
standard error $\sqrt{\langle N_3\rangle/\text{trials}}$, consistent
with the near-Poisson fluctuation of the count.

\subsection*{A.3 Perturbation pipeline}
For each prime $p$: (i) diagonalise $U_A$, cluster eigenvalues at
tolerance $10^{-4}$ on rounded real/imaginary parts; (ii) draw
$E\sim\mathrm{GUE}(N)$ and realise $e^{i\varepsilon E}$ through its
own eigendecomposition; (iii) match perturbed to unperturbed levels by
the Hungarian algorithm on the cost matrix
$C_{ij}=|\lambda_i^{\mathrm{new}}-\lambda_j^{\mathrm{old}}|$; the
maximum assignment distance scales linearly in $\varepsilon$
($7.6\times10^{-8}$ at $\varepsilon=10^{-5}$,
$7.6\times10^{-7}$ at $10^{-4}$), confirming first-order tracking;
(iv) group matched levels by their origin multiplet, take the three
largest multiplets, unfold intra-multiplet gaps to unit mean spacing,
pool them within the prime; (v) evaluate KS once per prime against
the GOE Wigner surmise CDF $1-e^{-\pi s^2/4}$ and against Poisson
$1-e^{-s}$; no cross-prime pooling is performed. Cluster sizes were
stable across $\varepsilon\in[10^{-5},10^{-3}]$ (e.g.\ $69\times5$ at
$p=2207$).

\subsection*{A.4 Reproducibility}
All scripts live in the repository with fixed seeds recorded in their
artefact headers; the artefact manifest maps every number above to
script and command. Compiled artefacts:
\texttt{unitary\_perturbation.json},
\texttt{run46\_exact\_mean.json}, \texttt{run36\_audit.json}.

\begin{thebibliography}{9}

\bibitem{BGS84}
O. Bohigas, M.-J. Giannoni, C. Schmit,
\emph{Characterization of chaotic quantum spectra and universality of
level fluctuation laws},
Phys.\ Rev.\ Lett.\ \textbf{52} (1984), 1--4.

\bibitem{HB80}
J.~H. Hannay, M.~V. Berry,
\emph{Quantization of linear maps on a torus --- Fresnel diffraction
by a periodic grating},
Physica D \textbf{1} (1980), 267--290.

\bibitem{KR00}
P. Kurlberg, Z. Rudnick,
\emph{Hecke theory and equidistribution for the quantization of linear
maps of the torus},
Duke Math.\ J.\ \textbf{103} (2000), no.~1, 47--77.

\bibitem{KR01}
P. Kurlberg, Z. Rudnick,
\emph{On quantum ergodicity for linear maps of the torus},
Comm.\ Math.\ Phys.\ \textbf{222} (2001), 201--227.

\bibitem{FNDB03}
F. Faure, S. Nonnenmacher, S. De Bi\`evre,
\emph{Scarred eigenstates for quantum cat maps of minimal periods},
Comm.\ Math.\ Phys.\ \textbf{239} (2003), 449--492.

\bibitem{OSV21}
A. Ostafe, I.~E. Shparlinski, J.~F. Voloch,
\emph{Equations and character sums with matrix powers, Kloosterman sums
over small subgroups and quantum ergodicity},
preprint, arXiv:2110.10941v2 [math.NT] (2021).

\end{thebibliography}

\end{document}
